%%%%%
%
\section{Sommation d'un param\`etre 1D}

\itemstyle{Op\'erateur~:} {\sf Som\_1D}
\index{Som\_1D}

%%%%%OBJET%%%%
\itemstyle{Objet~:}
Sommation de N \'el\'ements cons\'ecutifs 
dans le temps d'un param\`etre 1D.


%%%%%ARGUMENTS%%%%
\bigskip
%%%%%ENTREES%%%%
\begin{tabarg}
\multicolumn{3}{l}{\titretab {Entr\'ees}}\\
\hline
$X$ & Vecteur, {\tt Int16[$D$]} & Param\`etre d'entr\'ee \\
$D$ & Coeff., {\tt Int} & Nombre d'\'el\'ements d'un \'echantillon de $X$\\
$N$ & Coeff., {\tt Int} & Nombre de valeurs de la sommation\\
$Decal$ & Coeff., {\tt Int} & D\'ecalage (\textit{nombre de points})\\
\end{tabarg}

\medskip
%%%%%SORTIES%%%%
\begin{tabarg}
\multicolumn{3}{l}{\titretab {Sorties}}\\
\hline
$Xs$ & Vecteur, {\tt Int16[$D$]} & Param\`etre r\'esultat \\
\end{tabarg}


%%%%%FORMULE%%%%
\itemstyle{Formule}

%%Som
Pour chaque classe $k$ du param\`etre X : \\
\[
Xs_k[i] = \sum_{j=N\times i+Decal}^{N\times (i+1) -1 + Decal } X_k[j]
 \qquad i=0, Nbre\_echant(X)/N - 1
\]


%%%% COMMENTAIRE %%%%
\itemstyle{Commentaire~: }

Le nombre d'\'echantillons du vecteur r\'esultat est \'egal au nombre
d'\'echantillons du vecteur d'origine, divis\'e par  $N$, arrondi \`a
l'entier sup\'erieur.
Si le nombre d'\'echantillons de X n'est pas un multiple de N, la
derni\`ere valeur de $Xs$ devrait \^etre "valeur manquante".\\
La fr\'equence d'\'echantillonnage du param\`etre r\'esultat est le rapport de
la fr\'equence du param\`etre d'entr\'ee par $N$.

%%%%%AUTEUR%%%%
\itemstyle{Auteur~:}
CNRM/GMEI/TRAMM

%%%%%REFERENCE%%%%
\itemstyle{R\'ef\'erences~:}
